\chapter{الضرب: الخطوات الأولى}\label{ch:g2:mult}

أربع علب في كل واحدة ثلاث كعكات: بدل $3 + 3 + 3 + 3$ نكتب
$4 \times 3$. الضرب اختصار لجمع العدد نفسه مرات كثيرة --- وأحسن
صورة له شبكة مرتّبة من صفوف
وأعمدة.

\section{جمع العدد نفسه}

\begin{definition}[الضرب]\label{def:g2:mult:def}
$4 \times 3$ يعني «أربع مرات ثلاثة»:
\[
4 \times 3 = 3 + 3 + 3 + 3 = 12 .
\]
والإشارة هي $\times$، وتُقرأ «في». والناتج هو
\emph{الجداء}\index{جداء} (والضرب يكبر في \cref{ch:g3:mult}).
\end{definition}

\begin{omfigure}
\begin{tikzpicture}[scale=0.5]
  \foreach \r in {0,1,2,3}
    \foreach \c in {0,1,2}
      \fill[omDef] (\c,-\r) circle (0.28);
  \node[right, font=\small] at (3.0,-1.5)
    {$4$ صفوف من $3$ نقط: $4 \times 3 = 12$};
\end{tikzpicture}

{\small \emph{شبكة مستطيلة}: $4$ صفوف من $3$. والعدّ بالصفوف يعطي
$4 \times 3$؛ والنقط نفسها بالأعمدة تعطي $3 \times 4$ --- إذن
الترتيب لا يهمّ.}
\end{omfigure}

\begin{example}[الترتيب لا يهمّ]\label{ex:g2:mult:order}
$4 \times 3$ و $3 \times 4$ يعدّان النقط $12$ نفسها، مرة
بالصفوف ومرة بالأعمدة. وإدارة الشبكة على جنبها تبرهن ذلك. فتعلّم
\omterm{def:g2:mult:def}{جداء} واحد يعطي توأمه مجانًا.
\end{example}

\section{العدّ بالقفز والجداول السهلة}

\begin{method}[بناء جدول بالعدّ بالقفز]\label{met:g2:mult:skip}
لإيجاد $5 \times 4$، عُدّ بالأربعات خمس مرات:
$4, 8, 12, 16, 20$. فالعدّ بالقفز على مستقيم الأعداد
\emph{هو} الجدول.
\begin{itemize}
  \item جدول $2$: $2, 4, 6, 8, 10, \dots$ (الأضعاف)؛
  \item وجدول $5$: $5, 10, 15, 20, 25, \dots$ (ينتهي بالرقم $5$ أو
        $0$)؛
  \item وجدول $10$: $10, 20, 30, \dots$ (يكفي أن تضيف صفرًا).
\end{itemize}
\end{method}

\begin{omfigure}
\begin{tikzpicture}[scale=0.6]
  \draw[-Stealth] (-0.4,0) -- (11.4,0);
  \foreach \i in {0,...,11} \draw (\i,-0.1) -- (\i,0.1);
  \foreach \i in {0,2,4,6,8,10}
    \node[below, font=\footnotesize] at (\i,-0.12) {\i};
  \foreach \i in {0,1,2,3,4}
    \draw[-Stealth, omDef, thick] (2*\i,0.28) .. controls
      (2*\i+1,0.85) .. (2*\i+2,0.28);
  \node[omDef, font=\small] at (5.5,1.1) {قفزات بمقدار $2$};
\end{tikzpicture}

{\small جدول $2$ قفزاتٍ على المستقيم: $2, 4, 6, 8, 10$. وخمس
قفزات بمقدار $2$ تبلغ $10$، إذن $5 \times 2 = 10$.}
\end{omfigure}

\begin{example}[الأضعاف والأنصاف]\label{ex:g2:mult:double}
الضرب في $2$ تضعيف: $2 \times 7 = 14$. والعكس، وهو التنصيف،
يُلغيه: فنصف $14$ هو $7$. والأضعاف المحفوظة عن ظهر قلب تجعل
جدول $2$ كله فوريًّا --- وتمهّد للقسمة
(\cref{ch:g3:division}).
\end{example}

\section{الضرب في المسائل}

\begin{example}\label{ex:g2:mult:problem}
«في العلبة $5$ بيضات. فكم بيضة في $4$ علب؟» أربع مجموعات من
خمسة: $4 \times 5 = 20$ (بالعدّ بالقفز: $5, 10, 15, 20$). فيها
$20$ بيضة.
\end{example}

\section{تمارين}

\begin{exercise}[$\star$]\label{exo:g2:mult:1}
اُكتب على صورة ضرب، ثم احسب: $2 + 2 + 2$؛ \quad
$5 + 5 + 5 + 5$؛ \quad $10 + 10 + 10$.
\end{exercise}

\begin{exercise}[$\star$]\label{exo:g2:mult:2}
اُرسم شبكة من $3$ صفوف فيها $5$ نقط. فأي ضرب تُظهر؟ وماذا لو
أدرتها على جنبها؟
\end{exercise}

\begin{exercise}[$\star$]\label{exo:g2:mult:3}
عُدّ بالقفز لتملأ جدول $2$ حتى $2 \times 6$؛ وجدول
$5$ حتى $5 \times 6$.
\end{exercise}

\begin{exercise}[$\star$]\label{exo:g2:mult:4}
احسب: $2 \times 7$؛ \quad $5 \times 3$؛ \quad $10 \times 4$؛
\quad $2 \times 9$؛ \quad $5 \times 6$.
\end{exercise}

\begin{exercise}[$\star$]\label{exo:g2:mult:5}
احسب الأضعاف: ضعف $6$؛ وضعف $8$؛ وضعف $10$.
ثم الأنصاف: نصف $12$؛ ونصف $20$.
\end{exercise}

\begin{exercise}[$\star$]\label{exo:g2:mult:6}
احسب جدول $3$ بالعدّ بالقفز: $3, 6, 9, ?, ?, ?$. وكم يساوي
$3 \times 5$؟
\end{exercise}

\begin{exercise}[$\star$]\label{exo:g2:mult:7}
«هناك $4$ طاولات حول كل واحدة $5$ كراسٍ. فكم كرسيًّا؟»
اُكتب عملية الضرب والجواب.
\end{exercise}

\begin{exercise}[$\star$]\label{exo:g2:mult:8}
للعنكبوت $8$ أرجل. فكم رجلًا عند $3$ عناكب؟ (عُدّ بالقفز
بمقدار $8$: $8, 16, \dots$)
\end{exercise}

\begin{exercise}[$\star$]\label{exo:g2:mult:9}
طريقتان لترتيب $12$ نقطة في صفوف متساوية: اُرسم $2$ من الصفوف فيها
$6$، ثم $3$ صفوف فيها $4$. واكتب عملية الضرب لكل منهما.
\end{exercise}

\begin{exercise}[$\star\star$]\label{exo:g2:mult:10}
تأتي الكعكات في علب من $5$. وتحتاج مريم إلى $18$ كعكة لحفلة.
فكم علبة عليها أن تشتري؟ (عُدّ بالقفز بمقدار $5$ حتى تتجاوز $18$.)
\end{exercise}
